% Complete proof for the original screening architecture, with explicitly
% repaired comparison, confidence calibration, and physical stopping rules.
% This file does not prove a claim about CP-UCB-AA or inherit the former G3 bound.
\section{A fixed-budget analysis of the screening-based UCB-AA procedure}
\label{app:ucbaa_repaired}

\subsection{Exact scope and executable conventions}

Fix a deterministic physical horizon $T\ge1$ and $0<\delta<1$.
There are finitely many available arms by every finite time, at least one arm
is available at time one, and arrivals are deterministic and permanent.
The arm streams are independent, and each consists of iid $[0,1]$-valued
rewards with mean $r_i$. The proof also applies conditionally on an independent exogenous
arrival schedule. No gap-stability or cumulative-jump assumption is imposed.
The learner observes the arrivals before selecting the arm at each physical
time. Newly arriving arms are buffered until the next algorithmic round.
All observations are retained, and an eliminated arm is not sampled again.

Set
\begin{equation}
 M_T=\left\lceil\frac12\log_2(T+1)\right\rceil,\qquad
 N=\frac{8T4^{M_T}}{\delta},\qquad \widetilde\Delta_p=2^{-p},
 \label{eq:ucbaa_calibration}
\end{equation}
and, for $0\le p\le M_T$, retain the original target and radius formulas
\begin{equation}
 n_p=\left\lceil
       \frac{2\log(N\widetilde\Delta_p^2)}{\widetilde\Delta_p^2}
       \right\rceil,\qquad
 b_p=\sqrt{\frac{\log(N\widetilde\Delta_p^2)}{2n_p}}.
 \label{eq:ucbaa_targets}
\end{equation}
The calibration in \eqref{eq:ucbaa_calibration} is a sufficient choice for
this theorem; the smaller empirical calibration in the submitted version
is not covered by this probability guarantee.

Write $s_m$ for the physical starting time of algorithmic round $m$, with
$s_0=1$. At $s_m$ admit every arm in the buffer, including arrivals at $s_m$.
The admitted batch is $A_m$, its cumulative union is $A(m)$, the incumbent
survivors are $C_{m-1}$, and the initial active set is
$B_m=C_{m-1}\cup A_m$. Put $C_{-1}=\varnothing$.
These are the round-based sets of the original UCB-AA procedure. Its
notation $K(m)$, when used, denotes a prescribed nondecreasing upper
bound on $|A(m)|$; it is different from the actual physical-time count
$K_t=|\mathcal A_t|$. This proof does not require $K(m)$ as an additional
input because the calibration controls at most $T$ observed streams.

The following conventions make the original two-stage architecture precise.
They are part of the analyzed procedure, not conclusions about unspecified
implementations.
\begin{enumerate}
 \item An elimination call at phase $p$ first samples every arm in its
 candidate set $S$ until its lifetime count reaches $n_p$. An arm already
 at or above that target receives no additional sample. After the entire
 call's sampling is complete, form
 \[
  \mathrm{LCB}_i(p)=\hat r_i(n_p)-b_p,\qquad
  \mathrm{UCB}_i(p)=\hat r_i(n_p)+b_p
 \]
 from the actual first $n_p$ observations of each reference arm.
 Remove simultaneously from $S$ exactly those arms for which
 \[
  \mathrm{UCB}_i(p)<\max_{j\in R}\mathrm{LCB}_j(p).
 \]
 In particular, the inequality is strict. Prefix statistics for incumbents
 are retained; an unobserved prefix is never substituted into a comparison.
 \item For $m\ge2$, run preliminary phases $p=0,\ldots,m-2$ on the
 surviving part of $A_m$, with reference set equal to that part together
 with $C_{m-1}$. Only new arms are removed in these calls. Then run formal
 phases $p=m-1,m$ on the merged survivors. For $m=1$, run only formal
 phases $0,1$; for $m=0$, run formal phase $0$.
 \item If a new batch is empty, all its preliminary phases are no-ops:
 the set entering formal phase $m-1$ is $C_{m-1}$. An empty surviving new
 batch is handled the same way. This specifies the missing intermediate
 set in the original pseudocode when preliminary screening is skipped.
 \item Sampling orders and ties are fixed by observed history and
 independent policy randomization. Every physical pull counts toward
 $T$. Stop immediately after pull $T$, even in the middle of an arm's
 target, an elimination call, or a round. No terminal round is appended
 to remove buffered arrivals from the regret ledger.
\end{enumerate}

Thus the alteration preserves buffering, screening against incumbent
evidence, and formal elimination. It repairs the statistical bookkeeping
and stopping convention; it does not replace UCB-AA by a reward-service
mixture or by CP-UCB-AA.

\subsection{Three distinct contributions to regret}

At a physical time $t\le T$, let $m(t)$ be its current algorithmic round.
Define the best initially active mean and the best admitted mean by
\[
 r_{*(m)}=\max_{i\in B_m}r_i,\qquad
 r_{\mathrm{adm}}(m)=\max_{i\in A(m)}r_i.
\]
For an admitted arm in round $m$, distinguish the original roundwise gap
$\Delta_i^{\mathrm{rd}}(m)=r_{*(m)}-r_i$ from the physical-time gap
$\Delta_i(t)=r^*(t)-r_i$, where
$r^*(t)=\max_{i\in\mathcal A_t}r_i$. The superscript marks the
algorithmic-round clock; it does not alter the physical-time gap definition.
The following three truncated quantities have different meanings:
\begin{align}
 R_{\mathrm{round}}(T)
  &=\sum_{t=1}^T\{r_{*(m(t))}-r_{I_t}\},\nonumber\\
 \mathcal H_T
  &=\sum_{t=1}^T\{r_{\mathrm{adm}}(m(t))-r_{*(m(t))}\},\nonumber\\
 D_T
  &=\sum_{t=1}^T\{r^*(t)-r_{\mathrm{adm}}(m(t))\}.
 \label{eq:ucbaa_three_terms}
\end{align}
Here $\mathcal H_T$ is the deficit caused by previously eliminated admitted arms,
whereas $D_T$ is arrival delay alone. On every sample path,
\begin{equation}
 R_T(\pi)=R_{\mathrm{round}}(T)+\mathcal H_T+D_T.
 \label{eq:ucbaa_exact_identity}
\end{equation}
All three terms are nonnegative, since the selected arm belongs to $B_{m(t)}$
and $B_{m(t)}\subseteq A(m(t))\subseteq\mathcal A_t$.
In particular, the residual after subtracting roundwise regret is not in
general arrival delay alone.

For a prescribed complete run of $M+1$ algorithmic rounds, the original
$T_0(\pi)$ denotes its actual, generally random number of pulls, and
$R_{\mathrm{round}}(T_0(\pi))=R_M(\pi)$. The identity
\eqref{eq:ucbaa_exact_identity} applies at that completed endpoint as a
pathwise accounting statement. The fixed-budget theorem below does not
replace a deterministic $T$ by $\mathbb E T_0(\pi)$ or assert an
unproved bound at an arbitrary stopping time.

Arrival delay has an exact exposure representation, including an unfinished
last round. Let $\beta_u=r^*(u)-r^*(u-1)\ge0$ for $u\ge2$, and define
\[
 \alpha_T(u)=\min\bigl(\{s_m:u\le s_m\le T\}\cup\{T+1\}\bigr).
\]
Because every available arm is admitted at a round boundary,
$A(m)=\mathcal A_{s_m}$ and $r_{\mathrm{adm}}(m)=r^*(s_m)$.
Changing the order of two finite sums yields
\begin{equation}
 D_T=\sum_{u=2}^T\beta_u\{\alpha_T(u)-u\}.
 \label{eq:ucbaa_delay_exposure}
\end{equation}
An improvement arriving exactly at a round boundary has zero delay.
An improvement that is still buffered at $T$ is charged through $T$.

\begin{theorem}[UCB-AA at a fixed physical budget]
\label{thm:ucbaa_fixed_budget}
Under the model and executable conventions above, let
$K_T=|\mathcal A_T|$. With probability at least $1-\delta$, the best
admitted mean is always represented in the active set, $\mathcal H_T=0$, and
\begin{equation}
 R_T(\pi)\le
 \min\left\{T,\ 8\sqrt{(2\log N+1)K_TT}+D_T\right\}.
 \label{eq:ucbaa_fixed_hp}
\end{equation}
Consequently
\begin{equation}
 \mathbb E R_T(\pi)\le
 \min\left\{T,\ 8\sqrt{(2\log N+1)K_TT}
                   +\mathbb E D_T+\delta T\right\}.
 \label{eq:ucbaa_fixed_expectation}
\end{equation}
The confidence event and both bounds include the last incomplete round and
all samples in incomplete calls. The statement does not require a unique
best arm and does not compress an elimination-time gap into a terminal gap.
\end{theorem}

\begin{proof}
First check the targets and the stopping convention. Since
$\log(N4^{-M_T})=\log(8T/\delta)>1$, the function
$x\mapsto2x\log(N/x)$ is increasing on $[1,4^{M_T}]$.
More explicitly, its successive phase differences are
\[
 2\,4^{p-1}\{3\log(N4^{-p})-\log4\}>1,\qquad p\ge1.
\]
Therefore $n_p$ is strictly increasing. Also
\begin{equation}
 b_p\le\widetilde\Delta_p/2,\qquad
 n_p\le(2\log N+1)4^p,\qquad n_{M_T}>T.
 \label{eq:ucbaa_target_facts}
\end{equation}
The last fact follows from $4^{M_T}\ge T+1$ and
$\log(N4^{-M_T})>1$. Formal elimination with a strict inequality never
removes every arm, even off the confidence event: an arm with maximal lower
confidence bound cannot satisfy the removal rule because its upper bound
is at least its own lower bound. Initial and subsequent active sets are
therefore nonempty. Completing round $M_T$ would require a surviving arm
to have $n_{M_T}>T$ observations, and hence cannot happen before the
physical stopping time. The finite phase range suffices on every path.

For the global event, order arms by their first observed reward. At most
$T$ streams are exposed by time $T$. An arm's unobserved reward stream is
independent of the history that first selects it, and the bounded iid
Hoeffding inequality is uniform over its mean. For each exposed stream,
the union of its potential-prefix failures over $p=0,\ldots,M_T$ has
probability at most
\[
 \sum_{p=0}^{M_T}2\exp(-2n_pb_p^2)
  =\frac2N\sum_{p=0}^{M_T}4^p
  <\frac{8\,4^{M_T}}{3N}.
\]
Here the potential prefix is defined even if the policy never observes
that many rewards. Applying the same bound conditionally at the first
exposure and summing over at most $T$ exposed streams gives an event $E$
with
\begin{equation}
 \mathbb P(E^c)\le\frac{8T4^{M_T}}{3N}=\frac\delta3\le\delta,
 \qquad
 |\hat r_i(n_p)-r_i|\le b_p
 \label{eq:ucbaa_global_coverage}
\end{equation}
for every observed arm and every phase needed in a comparison.
One direct proof of the concentration inequality used here is as follows.
For a random variable $Z\in[0,1]$, the second derivative of its centered
log moment generating function is a tilted variance and is at most $1/4$.
Twice integrating gives
$\mathbb E e^{\lambda(Z-\mathbb EZ)}\le e^{\lambda^2/8}$.
Multiplication over independent observations and exponential Markov
inequality give the two-sided bound $2e^{-2nx^2}$.
The exposure argument uses only independence of an unrevealed stream,
not independence of its eventual sample count or admission round.
More formally, let $F_i$ be the union of potential-prefix failures for
stream $i$, and let $q=(2/N)\sum_{p=0}^{M_T}4^p$. Before revealing the
reward at time $t$, condition on the observed history and the policy's
current independent randomization. On the event that $I_t$ has never
been sampled, its stream is still unrevealed and the conditional
probability of $F_{I_t}$ is at most $q$. Summing over $t\le T$ bounds
the probability that any first-exposed stream fails by $Tq$.
This conditioning occurs before that stream's first reward is revealed.

Work on $E$. At every comparison, all reference means lie between their
actual lower and upper confidence bounds. A maximal-mean arm in the
reference set cannot satisfy the strict removal inequality. For formal
elimination this preserves the maximum of the entire active set. In
pre-elimination an incumbent maximum is not eligible for removal, and a
new maximum cannot satisfy the removal inequality either. Induction over
calls and rounds shows
\[
 \max_{i\in C_m}r_i=\max_{i\in A(m)}r_i
\]
at every completed round, and the best admitted mean is preserved during
every incomplete round as well. Thus
$r_{*(m)}=r_{\mathrm{adm}}(m)$ and $\mathcal H_T=0$.

There is no unavailable-prefix comparison in this induction. At the end
of round $m-1$, every incumbent survivor has exactly $n_{m-1}$ lifetime
observations. In preliminary phase $p\le m-2$, this includes the stored
first $n_p$ observations. New reference arms have just been sampled to
$n_p$ before the simultaneous comparison. Formal calls likewise have
the requisite prefixes after their sampling step.

Next establish the actual-count invariant. After any completed phase $p$
in round $m$, every eligible arm that survives the comparison has
\begin{equation}
 r_{*(m)}-r_i\le4b_p.
 \label{eq:ucbaa_survival_gap}
\end{equation}
Indeed a maximal-mean reference arm $j$ is still present, so
$\mathrm{LCB}_j(p)\ge r_{*(m)}-2b_p$, whereas
$\mathrm{UCB}_i(p)\le r_i+2b_p$. Survival requires the latter quantity
to be at least the former. This proves
\eqref{eq:ucbaa_survival_gap} without conditioning on a selected estimate.
The claim concerns eligible new arms in a preliminary call and all
survivors in a formal call; no such bound is asserted for untested
incumbents during preliminary screening.

Consider a physical pull that is the $n$th lifetime observation of its arm
and occurs while filling target $n_p$ in round $m$. If $p=0$, its roundwise
gap is at most one. If $p\ge1$ and the arm is new in the current round,
it survived phase $p-1$ of this same round before this pull, and its gap
is at most $4b_{p-1}$. If it is an incumbent, no sample is needed in
formal phase $m-1$, since it already has $n_{m-1}$ observations.
Its new samples occur only in phase $m$, after it survived the current
round's formal phase $m-1$, giving the same bound. This includes $m=1$;
in round zero the only target is $n_0$.
Consequently \eqref{eq:ucbaa_target_facts} gives, for every actual pull,
\begin{equation}
 r_{*(m)}-r_i\le
 \min\left\{1,\ 4\sqrt{\frac{2\log N+1}{n}}\right\}.
 \label{eq:ucbaa_actual_prefix_charge}
\end{equation}
For $p\ge1$, use $4b_{p-1}\le4\widetilde\Delta_p$ and
$n\le n_p\le(2\log N+1)4^p$. For $p=0$, use
$n\le n_0\le2\log N+1$. A partially sampled final phase requires
only the preceding comparison, which has already finished, so the
argument does not presume completion of the final target.

No arm's lifetime count is reset. Summing the distinct charges gives
\begin{align*}
 R_{\mathrm{round}}(T)
 &\le4\sqrt{2\log N+1}
       \sum_{i:N_i(T)>0}\sum_{n=1}^{N_i(T)}n^{-1/2}\\
 &\le8\sqrt{2\log N+1}\sum_i\sqrt{N_i(T)}
 \le8\sqrt{(2\log N+1)K_TT}.
\end{align*}
The second inequality follows by integrating $x^{-1/2}$ on each unit
interval; the last is Cauchy--Schwarz and $\sum_iN_i(T)=T$.
Combine this with \eqref{eq:ucbaa_exact_identity}, $\mathcal H_T=0$, and the
pathwise bound $R_T\le T$ to obtain \eqref{eq:ucbaa_fixed_hp}.
Finally use $R_T\le T$ on $E^c$, nonnegativity of $D_T$, and
\eqref{eq:ucbaa_global_coverage} to obtain
\eqref{eq:ucbaa_fixed_expectation}. The expectation does not condition
the law of selected rewards on coverage.
\end{proof}

For example, choosing $\delta=T^{-2}$ for $T\ge2$ gives
$\log N=O(\log(T+1))$. Thus the fixed-budget bound is sublinear whenever
$K_T\log(T+1)=o(T)$ and $\mathbb E D_T=o(T)$.
This is a sufficient arrival-delay condition stated on the physical
timeline, not a guarantee of sublinear regret for arbitrary arrivals.
The theorem is a corrected upper bound; no matching minimax claim follows
merely from the square-root form of its statistical term.

\Needspace{12\baselineskip}
\subsection{Static arms and fixed-gap arrivals}

\begin{corollary}[An initially available global optimum]
\label{cor:ucbaa_fixed_gap}
Under Theorem~\ref{thm:ucbaa_fixed_budget}, suppose an arm available at
time one has mean $r^\star$, and no arm arriving by $T$ has mean larger
than $r^\star$. Put $\Delta_i=r^\star-r_i$ and $A=2\log N+1$.
Then $D_T=0$ on every sample path. With probability at least $1-\delta$,
\begin{equation}
 R_T(\pi)\le
 \min\left\{T,\ \sum_{\substack{i\in\mathcal A_T\\\Delta_i>0}}
   \min\left\{\Delta_i(T-\tau_i+1),\frac{16A}{\Delta_i}\right\}
     \right\}.
 \label{eq:ucbaa_fixed_gap_bound}
\end{equation}
The same deterministic sum bounds $\mathbb E R_T(\pi)$ after adding
$\delta T$ and taking the minimum with $T$.
\end{corollary}
\begin{proof}
The physical benchmark equals $r^\star$ throughout the horizon. Because
an optimal arm is in the initial admitted batch, the best admitted mean
also equals $r^\star$ on every path, even if that arm is later eliminated.
Thus $D_T=0$ pathwise. On the common coverage event, the active benchmark
is $r^\star$ as well, so each pull's roundwise gap in
\eqref{eq:ucbaa_actual_prefix_charge} is its fixed physical gap $\Delta_i$.
If a suboptimal arm has $N_i(T)>0$ observations, apply that inequality to
its last observed prefix $n=N_i(T)$ to obtain
\[
 N_i(T)\le\frac{16A}{\Delta_i^2}.
\]
This is an inequality for an integer count, so it requires no extra
rounding term. It holds trivially at zero count. Independently,
$N_i(T)\le T-\tau_i+1$ by one pull per physical time and no prearrival
sampling. The exact fixed-gap identity
$R_T(\pi)=\sum_i\Delta_iN_i(T)$ gives
\eqref{eq:ucbaa_fixed_gap_bound}; arms tied at the optimum contribute
zero. On the exceptional event use $R_T\le T$.
\end{proof}

For $\delta=T^{-2}$ and $T\ge2$, $A=O(\log(T+1))$.
For a static finite arm set, this recovers the familiar upper order
$O\{\sum_{i:\Delta_i>0}\log(T+1)/\Delta_i\}$.
For arriving distractors satisfying $\Delta_i\ge g>0$, it gives
$O\{K_T\log(T+1)/g\}$, and hence $O(K_T\log(T+1))$ for fixed $g$.
These are special-case upper bounds, not an assertion that preliminary
screening is necessary for this order or a comparison proving linear
regret for a procedure without preliminary screening.

The original fixed-gap illustration can be expressed on the exogenous
physical clock: release an optimal arm of mean $3/4$ and an arm of mean
$1/4$ at time one, and at deterministic physical times $\tau_i$ release
arms of means $r_i=1/2-2^{-i-2}$, $i\ge1$. Their gaps are
$1/4+2^{-i-2}\ge1/4$, so the bound is $O(K_T\log(T+1))$.
The release times here are external physical times, not algorithmic round
indices chosen after seeing the procedure's random duration.

\subsection{A physical-time doubling wrapper}

For completeness an unknown horizon can be handled by restarting at
physical times $1,2,4,\ldots$. Epoch $\ell\ge0$ has prescribed length
$h_\ell=2^\ell$ and local failure budget
$\delta_\ell=\delta/(2h_\ell^2)$. At its start, initialize UCB-AA with
all currently available arms as its initial batch, empty local counts,
and the calibration \eqref{eq:ucbaa_calibration} for horizon $h_\ell$
and failure budget $\delta_\ell$. Previously eliminated arms re-enter.
Use only new rewards for this epoch's prefix statistics. Within the epoch,
retain the two original stages and buffer new arrivals until the next local
round. Stop the last epoch at the actual physical endpoint $T$.
This is a standard physical-budget restart, distinct from completing
geometrically many algorithmic rounds regardless of their random duration.

Let $D_T^{\mathrm{dbl}}$ sum the arrival-delay terms of these epochs,
with each epoch's unfinished local round truncated at its epoch endpoint
or at $T$. Put
\[
 L_T=2\log(128T^4/\delta)+1.
\]
\begin{corollary}[A valid unknown-horizon extension]
\label{cor:ucbaa_physical_doubling}
For every deterministic physical endpoint $T\ge1$, the wrapper satisfies
\[
 \mathbb P\left\{R_T(\pi)\le
       32\sqrt{L_TK_TT}+D_T^{\mathrm{dbl}}\right\}\ge1-\delta
\]
and
\[
 \mathbb E R_T(\pi)\le
 \min\left\{T,\ 32\sqrt{L_TK_TT}
                  +\mathbb E D_T^{\mathrm{dbl}}+\delta\right\}.
\]
The high-probability inequality can be made simultaneous over all physical
endpoints on one event of probability at least $1-\delta$.
\end{corollary}
\begin{proof}
Conditional on the history at an epoch start, future unrevealed rewards
are independent arm streams, each consisting of iid rewards with its
original mean. This follows by
successively exposing the next unrevealed observation of a selected arm;
selection and the epoch boundary depend only on already revealed entries.
The fixed-budget confidence argument therefore applies conditionally,
and its actual-prefix proof holds for every truncation of that epoch.
The union of epoch failures has probability at most
$\sum_{\ell\ge0}\delta_\ell=2\delta/3<\delta$.

If $M_\ell=\lceil\tfrac12\log_2(h_\ell+1)\rceil$, then
$4^{M_\ell}\le4(h_\ell+1)\le8h_\ell$ and
\[
 N_\ell=\frac{8h_\ell4^{M_\ell}}{\delta_\ell}
       \le\frac{128h_\ell^4}{\delta}.
\]
For every epoch that starts by $T$, its local $2\log N_\ell+1$ is at
most $L_T$. On the common good event, sum the fixed-budget bounds over
complete epochs and the last partial epoch. Their available-arm counts
are at most $K_T$. Since
\[
 \sum_{\ell=0}^{\lfloor\log_2 T\rfloor}\sqrt{2^\ell}
 \le(2+\sqrt2)\sqrt T<4\sqrt T,
\]
their statistical terms sum to at most $32\sqrt{L_TK_TT}$.
This proves the simultaneous high-probability statement.

For the expectation apply the good-event bound separately within each
epoch and charge at most $h_\ell$ regret for that epoch's failure.
Earlier erroneous eliminations do not propagate through the restart.
The total exceptional-event contribution is at most
\[
 \sum_{\ell\ge0}h_\ell\delta_\ell
    =\frac\delta2\sum_{\ell\ge0}2^{-\ell}=\delta.
\]
Add the nonnegative expected delay terms and use $R_T\le T$.
\end{proof}

\paragraph{Audit and boundaries of the repair.}
The proof uses only bounded-stream concentration, strict interval
comparisons, the roundwise survivor-count invariant, a finite sum, and
Cauchy--Schwarz; each dependency is proved above. The competing route of
terminal-gap compression is unnecessary and is not reinstated. Ties,
empty new batches, the first two rounds, singleton active sets,
deterministic rewards, arbitrary within-round benchmark jumps, and an
unfinished first or last target are all covered. The historical-deficit
term vanishes only on the common confidence event. The delay term is
retained explicitly, including at the endpoint. The unknown-horizon
wrapper restarts the original screening algorithm and does not claim
history retention across epochs. Neither result establishes that the
original smaller empirical calibration has the same error probability,
that the withdrawn gap-dependent corollary is true, or that the old
round-based online theorem has been repaired without changing its clock.
